\begin{sidewaystable}[p]
\centering
\scriptsize
\caption{Контрасты первичного исхода O1 по четырём процедурам, оба estimand и оба контраста}
\label{tab:3}
\begin{tabular}{lllllllll}
\toprule
Контраст & Estimand & Процедура & Популяция контраста & Пар & Эффект, сырая шкала & 95\% CI & p & Знак после конвенции \\
\midrule
P3 − P1 & full & 1 — индекс стиля score-v2 & весь корпус, 45 заданий & 359 & −1.7655 & [−2.6639; −0.8773] & 0.0004 & − \\
P3 − P1 & full & 2 — классификатор P2b & только жанр seo, 15 заданий & 120 & −0.0941 & [−0.1473; −0.0463] & 0.0028 & − \\
P3 − P1 & full & 3 — zero-shot nll-v2 & весь корпус, 45 заданий & 359 & +0.1140 & [0.0964; 0.1313] & 0.0004 & − \\
P3 − P1 & full & 4 — модель-судья judge-v2 & весь корпус, 45 заданий & 359 & +0.5153 & [0.1806; 0.8552] & 0.0032 & + \\
P3 − P1 & net & 1 — индекс стиля score-v2 & весь корпус, 45 заданий & 359 & +2.0264 & [0.8953; 3.1429] & 0.0016 & + \\
P3 − P1 & net & 2 — классификатор P2b & только жанр seo, 15 заданий & 120 & +0.0169 & [−0.0134; 0.0496] & 0.3168 & + \\
P3 − P1 & net & 3 — zero-shot nll-v2 & весь корпус, 45 заданий & — & не определён & — & — & — \\
P3 − P1 & net & 4 — модель-судья judge-v2 & весь корпус, 45 заданий & — & не определён & — & — & — \\
P2 − P1 & full & 1 — индекс стиля score-v2 & весь корпус, 45 заданий & 360 & +3.2812 & [2.0186; 4.6068] & 0.0004 & + \\
P2 − P1 & full & 2 — классификатор P2b & только жанр seo, 15 заданий & 120 & +0.0523 & [0.0325; 0.0738] & 0.0002 & + \\
P2 − P1 & full & 3 — zero-shot nll-v2 & весь корпус, 45 заданий & 360 & +0.0125 & [−0.0065; 0.0298] & 0.1988 & − \\
P2 − P1 & full & 4 — модель-судья judge-v2 & весь корпус, 45 заданий & 360 & +0.5125 & [0.0181; 1.0472] & 0.0392 & + \\
P2 − P1 & net & 1 — индекс стиля score-v2 & весь корпус, 45 заданий & 360 & −1.6860 & [−3.2913; −0.0343] & 0.0416 & − \\
P2 − P1 & net & 2 — классификатор P2b & только жанр seo, 15 заданий & 120 & +0.0722 & [0.0395; 0.1036] & 0.0014 & + \\
P2 − P1 & net & 3 — zero-shot nll-v2 & весь корпус, 45 заданий & — & не определён & — & — & — \\
P2 − P1 & net & 4 — модель-судья judge-v2 & весь корпус, 45 заданий & — & не определён & — & — & — \\
\bottomrule
\end{tabular}
\begin{minipage}{\linewidth}\footnotesize\vspace{2pt}
\textit{Примечание.} Единица анализа — пара текстов, порождённых одной моделью по одному заданию в двух режимах. Знаменатель: 359 или 360 пар у процедур 1, 3 и 4 при 45 кластерах; 120 пар у процедуры 2 при 15 кластерах, поскольку она считает контраст только на жанре seo. Данные: серия v2, статус current. Интервалы 95\%, кластерный бутстрап по заданию; кластер — задание, заданий 45. Публикуется абсолютная разница со своим интервалом; единого числа по четырём процедурам нет, сводный эффект не рассчитывался. Шкала: Конвенция интерпретации: больше значит более AI-подобно. У процедуры 3 она инвертирует знак сырой метрики, поэтому величины между процедурами не сравниваются. Сокращения: estimand full — полный балл, net — балл без структурных компонентов; P1, P2 и P3 — режимы задания. Пропуски: estimand net определён только у процедур 1 и 2; у zero-shot NLL и у судьи разложения на компоненты нет по построению, и ячейка означает «величина не определена», а не «не посчитана». Семейство первичных исходов закрыто поправкой Бонферрони на два зарегистрированных теста, α = 0.025 на тест; интервалы нескорректированные.
\end{minipage}
\end{sidewaystable}
